%% ============================================================
%%  SECCIÓN 6 — Procedimiento Experimental
%% ============================================================

\section{Procedimiento Experimental}
\label{sec:procedimiento}

El desarrollo experimental en el laboratorio de Física III se llevó a cabo siguiendo una secuencia metrológica rigurosa y ordenada, estructurada en los siguientes pasos cronológicos:

\begin{enumerate}
    \item \textbf{Encendido y calibración del generador de funciones (Pasos a--d):} Encendimos el generador de funciones OWON AG 051F y seleccionamos como forma de onda la señal periódica cuadrada (\texttt{Square}). Ajustamos la frecuencia de salida exactamente en $f = \qty{250}{\hertz}$, lo que corresponde a un período total de $T = \qty{4,00}{\milli\second}$ y un semiperíodo de conmutación de $T/2 = \qty{2,00}{\milli\second}$. Configuramos la amplitud pico a pico en \qty{10,0}{\voltpp} con un ciclo de trabajo (\texttt{Duty Cycle}) del \qty{50}{\percent}, asegurando la simetría entre el lapso de nivel alto (\qty{10}{\volt}) y el lapso de nivel bajo (\qty{0}{\volt}).
    
    \item \textbf{Configuración inicial del osciloscopio (Pasos e--h):} Encendimos el osciloscopio digital Rohde \& Schwarz HMO1202 y conectamos la salida del generador al Canal 1 (CH1) mediante un cable coaxial BNC. Ajustamos la perilla de la base de tiempo horizontal en \qty{500}{\micro\second\per\division}, logrando que un ciclo completo de \qty{4,0}{\milli\second} abarque con precisión 8 divisiones horizontales de la retícula de la pantalla. Configuramos la escala de ganancia vertical en \qty{2,0}{\volt\per\division}, de modo que la amplitud total de la señal cuadrada (\qty{10,0}{\volt}) ocupe exactamente 5 divisiones verticales. Esta calibración fue fundamental para nuestras lecturas: al valer cada división vertical \qty{2,0}{\volt} (o \qty{20}{\percent} de la señal completa), el nivel teórico del \qty{63,2}{\percent} ($V_C = \qty{6,32}{\volt}$) se localiza de forma inmediata a $3{,}16$ divisiones verticales sobre la línea de tierra de referencia.
    
    \item \textbf{Montaje del circuito RC serie con tierra común (Pasos i--k):} Desconectamos el cable coaxial directo e implementamos sobre el módulo acrílico el circuito serie formado por la resistencia $R_1 = \qty{9,98}{\kilo\ohm}$ y el condensador $C_1 = \qty{10,0}{\nano\farad}$ (Figura 6 de la guía del laboratorio). Para monitorizar ambos componentes de forma simultánea sin provocar cortocircuitos entre las masas de los canales, conectamos los terminales de tierra (cocodrilos negros) de las sondas de CH1 y CH2 en el nodo común central situado entre el resistor y el condensador. El terminal activo de la sonda de CH1 se conectó al extremo libre del condensador para registrar la tensión $V_C(t)$, mientras que el terminal activo de CH2 se conectó al extremo de la resistencia alimentado por el generador. Dado que la referencia común se ubicó en el nodo medio, invertimos la polaridad lógica del Canal 2 en el menú del osciloscopio (\texttt{CH2 -> INVERT}) para visualizar la caída $V_R(t) = R I(t)$ con polaridad positiva durante la fase de carga.
    
    \item \textbf{Adquisición de la constante de tiempo $\tau$ mediante cursores digitales (Pasos l--n):} Con las formas de onda estabilizadas en pantalla mediante disparo por flanco ascendente (\texttt{Edge Trigger}), activamos el menú de cursores temporales (\texttt{CURSOR -> TIME}). Posicionamos el cursor horizontal de referencia en el origen temporal del flanco ascendente ($t = 0$, $V = \qty{0}{\volt}$). A continuación, desplazamos el segundo cursor sobre la traza exponencial de CH1 ($V_C(t)$) hasta interceptar exactamente el valor de tensión correspondiente a $V_C = 0{,}632 \times \qty{10,0}{\volt} = \qty{6,32}{\volt}$ (morfología documentada en la Fig.~\ref{fig:carga_descarga}). El osciloscopio calculó de manera automática la diferencia temporal $\Delta t = t_2 - t_1$, la cual corresponde directamente a la constante de tiempo experimental $\tau_{\text{exp}}$. Repetimos la comprobación sobre la traza de CH2 ($V_R(t)$), ubicando el cursor en el instante en que la tensión decae hasta $V_R = 0{,}368 \times \qty{10,0}{\volt} = \qty{3,68}{\volt}$, verificando la coincidencia micrométrica entre ambos métodos.
    
    \item \textbf{Verificación en tiempo real de la Ley de Kirchhoff (Paso o):} Presionamos el botón frontal \texttt{MATH} del osciloscopio HMO1202 y seleccionamos la operación matemática de suma \texttt{ADD} entre los canales 1 y 2 ($\text{MATH} = \text{CH1} + \text{CH2}$). Observamos la aparición de una traza roja continua perfectamente horizontal en \qty{10,0}{\volt} durante el primer semiperíodo y en \qty{0}{\volt} durante el segundo semiperíodo (como se ilustra en la Fig.~\ref{fig:suma_kvl}), confirmando visual y cuantitativamente que la suma punto a punto $V_C(t) + V_R(t)$ anula las curvaturas exponenciales y reconstruye la señal cuadrada original del generador.
    
    \item \textbf{Variación sistemática de componentes (Pasos p--s):} Manteniendo inalterados la frecuencia (\qty{250}{\hertz}) y el voltaje de excitación (\qty{10,0}{\volt}), sustituimos sucesivamente el condensador por $C_2 = \qty{22,40}{\nano\farad}$ y $C_3 = \qty{33,23}{\nano\farad}$, registrando en cada caso el valor de $\tau$ en la hoja de datos de la Mesa C7. Posteriormente, reemplazamos la resistencia $R_1$ por los valores $R_2 = \qty{6,832}{\kilo\ohm}$ y $R_3 = \qty{2,250}{\kilo\ohm}$, repitiendo las mediciones para completar la matriz de 9 ensayos experimentales.
    
    \item \textbf{Montaje del circuito mixto con divisor resistivo (Paso 20 / Paso t):} Finalmente, implementamos el circuito de la Figura 9 y 10 de la guía de laboratorio, conectando la resistencia $R_1 = \qty{9,98}{\kilo\ohm}$ en serie con el bloque paralelo formado por la resistencia $R_2 = \qty{6,832}{\kilo\ohm}$ y el condensador $C_3 = \qty{33,23}{\nano\farad}$ (cuya respuesta se analiza en la Fig.~\ref{fig:paso20_thevenin}). Energizamos el circuito con la misma onda cuadrada de \qty{10,0}{\volt} y medimos con los cursores la tensión máxima de saturación en bornes del condensador $V_{C,\max}$, la constante de tiempo del transitorio resultante $\tau_{\text{th}}$, así como las corrientes pico y de régimen estacionario a través de la caída de potencial en $R_1$.
\end{enumerate}